% ============================================================
% !TeX root = ../main.tex
% Chapter 1 — Introduction
% ============================================================

\chapter{Introduction}

\label{ch:introduction}

\noindent
This chapter introduces the clinical and technical motivation for reconstructing
three-dimensional left-ventricular geometry from sparse cardiac magnetic
resonance imaging (MRI). It defines the reconstruction problem, presents the
research questions and objectives, and summarises the main contributions of
this thesis.

\section{Motivation}

\label{sec:motivation}

\noindent
Understanding cardiac structure is fundamentally a geometric problem. The
left ventricle (LV) is a continuously varying three-dimensional structure whose
shape, size, and wall thickness change throughout the cardiac cycle. These
properties provide important information about cardiac function and structural
remodelling, with changes in ventricular geometry and myocardial thickness
being associated with several forms of cardiac disease
\parencite{grossman1974wallthickness,demarvao2014atlas,medrano2014lvshapevariation}.

Cardiac magnetic resonance imaging (MRI) provides detailed observations of
this anatomy, but these observations are acquired as a set of two-dimensional
slices rather than as a complete three-dimensional representation. In a
short-axis (SAX) acquisition, the LV is sampled from the base towards the
apex, allowing the endocardial and epicardial boundaries to be delineated on
individual slices. Within each slice, the spatial resolution can be relatively
high, while the distance between adjacent slices is substantially larger.
Consequently, a structure that is anatomically continuous is represented by a
limited number of discrete observations.

This distinction becomes important when measurements are extended beyond the
acquired slices. Myocardial wall thickness, for example, is not simply a
property of isolated contours. It varies across the ventricle and can reflect
changes in cardiac structure, including myocardial hypertrophy
\parencite{grossman1974wallthickness}. Estimating such quantities from sparse
observations therefore depends not only on the accuracy of the segmentation,
but also on how well the underlying three-dimensional geometry is represented.

A continuous reconstruction can provide a more coherent description of the
LV and enable measurements in regions that are not directly observed. Such
representations are particularly relevant in large imaging studies, where
consistent anatomical descriptions are needed for automated analysis and
comparison across subjects
\parencite{bai2015biventricularatlas,bai2018automated}. They can also support
regional analysis of myocardial structure, including measurements organised
according to the American Heart Association 17-segment model
\parencite{cerqueira2002segmentation}.

The central challenge addressed in this thesis is therefore not the extraction
of LV contours alone, but the inference of the three-dimensional geometry that
lies between sparse observations. A useful reconstruction must recover
continuous endocardial and epicardial surfaces while preserving patient-specific
shape and maintaining a anatomically plausible myocardial wall. Simple
interpolation or surface lofting can provide continuity, but does not
necessarily recover the geometry that best represents the underlying anatomy.

This thesis investigates whether a phase-conditioned implicit neural
representation can learn this reconstruction problem from synthetic LV
geometries and meshes derived from real cardiac MRI segmentations. The proposed
approach represents the LV surfaces as continuous functions in three-dimensional
space, allowing geometry to be queried at arbitrary locations rather than only
at the acquired slices. The reconstructed surfaces are subsequently used to
study myocardial wall thickness and to compare the resulting measurements with
those obtained from RBF-derived clinical geometry. The aim is not to replace
clinical interpretation, but to investigate whether continuous geometric
representations can provide a more consistent basis for quantitative analysis
of sparse cardiac MRI.

\section{Problem Statement}

\label{sec:problem}

\noindent
SAX cardiac MRI provides detailed cross-sectional information about the LV, but
does not directly provide a complete three-dimensional representation. The
slices are separated along the ventricular long axis, leaving portions of the
geometry unobserved between adjacent acquisitions. As a result, the spatial
resolution within each slice is typically much finer than the resolution
between slices. Reconstructing a continuous surface from these sparse
observations is therefore an ill-posed geometric problem.

A simple connection of the segmented contours may produce a continuous surface,
but continuity alone does not guarantee a plausible reconstruction. Such
approaches can introduce irregularities, gaps, or local distortions, especially
when the distance between slices is large. The challenge is particularly
important when both the endocardial and epicardial surfaces are considered,
because they must be reconstructed as a consistent pair. The endocardium
defines the inner boundary of the LV cavity, while the epicardium defines the
outer boundary of the myocardium. Their relative position must remain
physiologically plausible so that the reconstructed wall retains positive and
spatially coherent thickness.

The problem is further constrained by the characteristics of publicly
available cardiac MRI datasets. These datasets commonly provide images and
segmentation masks for individual slices rather than complete three-dimensional
LV meshes. Meshes are nevertheless useful for directly training and evaluating
surface-reconstruction models because they provide an explicit representation
of the continuous geometry. A method is therefore required to bridge the gap
between sparse slice-based observations and continuous three-dimensional
surfaces.

This thesis addresses this problem using a phase-conditioned implicit neural
representation (INR) trained on synthetic LV shapes and meshes derived from
real cardiac MRI segmentations. The model receives sparse contours as input
and represents the endocardial and epicardial surfaces as continuous
three-dimensional functions. Conditioning the representation on cardiac phase
allows the model to distinguish between end-diastolic and end-systolic
geometry. In addition, the reconstruction is constrained so that the outer
surface remains outside the inner surface, preserving a positive myocardial
wall thickness.

The reconstructed geometry is evaluated at both the surface and measurement
levels. Geometric agreement is assessed against reference surfaces, while
myocardial wall thickness is measured on the reconstructed geometry and
compared with measurements obtained from RBF-derived clinical geometry. This
makes it possible to evaluate not only whether the proposed model produces
visually continuous surfaces, but also whether the resulting geometry preserves
quantities relevant to cardiac analysis.

\section{Research Questions}

\label{sec:research-questions}

\noindent
This thesis addresses the following research questions:

\begin{enumerate}[label=\textbf{RQ\arabic*.}]
\item To what extent can the proposed model reconstruct continuous,
watertight endocardial and epicardial LV surfaces from sparse SAX contours,
and does it improve geometric agreement over matched linear contour
lofting?

\item Can the phase-conditioned model represent end-diastolic and
end-systolic LV geometry while maintaining a positive myocardial-wall
separation between the endocardial and epicardial surfaces?

\item To what extent do local and regional wall-thickness measurements on
the reconstructed geometry agree with those obtained from
RBF-derived clinical geometry, and how does that agreement depend on the
measurement method?
\end{enumerate}
